\chapter{Photophobia (Light Sensitivity)}

\section*{Definition}
Photophobia is abnormal sensitivity to light, causing discomfort or pain when exposed to bright or even normal light levels. It is a neurological symptom involving connections between the trigeminal nerve and visual system.

\section*{What Happens in Your Body}
Photophobia involves intracranial photosensitivity mediated by melanopsin-expressing retinal ganglion cells that project to the trigeminal nerve system. In ocular photophobia, inflammation of the uveal tract, cornea, or meninges directly irritates trigeminal nerve endings. In neurological photophobia (e.g., migraine), the trigeminovascular system is sensitized, and light triggers or worsens headache through the retino-trigeminal pathway.

\section*{Causes}
\begin{itemize}
  \item Migraine (most common neurological cause -- light worsens headache)
  \item Meningitis (infectious or inflammatory irritation of meninges)
  \item Corneal conditions: Abrasion, keratitis, conjunctivitis, foreign body
  \item Uveitis (inflammation of the uveal tract)
  \item Iritis or iridocyclitis
  \item Post-operative (cataract surgery, refractive surgery)
  \item Concussion or traumatic brain injury
  \item Blepharospasm (involuntary eyelid closure)
  \item Medications: Antihistamines, anticholinergics, amphetamines
  \item Albinism (lack of retinal pigment reduces light absorption)
  \item Sleep deprivation (increased cortical sensitivity)
\end{itemize}

\section*{When to See a Doctor}
See your doctor if:
\begin{itemize}
  \item Photophobia with severe headache (suggests migraine or meningitis)
  \item Fever, neck stiffness, altered mental status (meningeal irritation -- emergency)
  \item Eye pain, redness, or visual changes (ocular cause)
  \item Photophobia following head trauma
  \item Associated nausea, vomiting, or visual aura (migraine)
  \item Chronic photophobia affecting daily function
\end{itemize}

\section*{Related Symptoms}
\begin{itemize}
  \item Headache
  \item Eye pain
  \item Red eye
  \item Blurred vision
  \item Nausea
  \item Visual aura
  \item Neck stiffness
  \item Blepharospasm
  \item Fatigue
\end{itemize}

\section*{Self-Care / Home Management}
\begin{itemize}
  \item Wear UV-protective sunglasses or photochromic lenses outdoors.
  \item Reduce screen brightness and enable night mode on devices.
  \item Use dim, indirect lighting at home; avoid fluorescent lighting.
  \item Wear a wide-brimmed hat in bright environments.
  \item Apply cool compresses to closed eyes for comfort.
  \item Rest in a dark, quiet room during acute episodes.
  \item Seek immediate care if photophobia is accompanied by fever, stiff neck, or severe headache.
\end{itemize}

\section*{How Your Doctor Will Evaluate You}
\begin{itemize}
  \item Slit-lamp eye examination to assess cornea, anterior chamber, and uvea.
  \item Visual acuity testing.
  \item Neurological examination including cranial nerve assessment.
  \item Lumbar puncture if meningitis suspected (fever, neck stiffness, altered mental status).
  \item MRI or CT brain if intracranial pathology suspected.
  \item Migraine diagnostic criteria assessment.
\end{itemize}

\section*{Treatment Options}
\begin{itemize}
  \item Treat underlying cause: Abortive/migraine prophylaxis for migraine, antibiotics for meningitis, cycloplegic drops for uveitis.
  \item Dark therapy and environmental light modification.
  \item Tinted lenses (FL-41 tint) for migraine-associated photophobia.
  \item Artificial tears and ocular surface protection for corneal causes.
  \item Botulinum toxin for blepharospasm.
  \item Cognitive behavioral therapy for chronic photophobia.
\end{itemize}

\section*{References}
\begin{thebibliography}{99}
\bibitem{photophobia1} Digre KB, Brennan KC. "Shedding light on photophobia." J Neuroophthalmol. 2012;32(1):68-81.
\bibitem{photophobia2} Katz BJ, Digre KB. "Diagnosis and management of photophobia." Semin Neurol. 2019;39(6):722-731.
\bibitem{photophobia3} Lebensohn JE. "Photophobia: mechanism and treatment." Am J Ophthalmol. 1951;34(9):1294-1300.
\bibitem{photophobia4} Custer PL, Gross CP, Trogdon J. "Photophobia." In: Cornea. 4th ed. Elsevier; 2017.
\bibitem{photophobia5} Noseda R, Copenhagen D, Burstein R. "Current understanding of photophobia." Cephalalgia. 2019;39(11):1479-1490.
\bibitem{photophobia6} Amin RM, Gokhale M, Vogel EA. "Photophobia in migraine: a symptom cluster." Curr Pain Headache Rep. 2020;24(10):58.
\end{thebibliography}

\clearpage
